\graphicspath{{notes/3canonical/asy/}}
\def\at#1#2{#1_{#2}}

\setcounter{section}{2}

\section{Canonical Forms}

\subsection{Jordan Forms \& Generalized Eigenvectors}

Throughout this course we've concerned ourselves with variations of a general question: for a given  map $\rT\in\cL(V)$, find a basis $\beta$ such that the matrix $[\rT]_\beta$ is as close to diagonal as possible. In this chapter we see what is possible when $\rT$ is non-diagonalizable.


\begin{example}{}{jordan1}
The matrix $A=\begin{smatrix}
-8&4\\-25&12
\end{smatrix}\in M_2(\R)$ has characteristic equation
\[p(t)=(-8-t)(12-t)+4\cdot 25=t^2-4t+4=(t-2)^2\]
and thus a single eigenvalue $\lambda=2$. It is non-diagonalizable since the eigenspace is one-dimensional
\[E_2 =\cN\begin{pmatrix}
-10&4\\-25&10
\end{pmatrix}=\Span\twovec 25\]
However, if we consider a basis $\beta=\{\vv_1,\vv_2\}$ where $\vv_1=\stwovec 25$ is an eigenvector, then $[\rL_A]_\beta$ is upper-triangular, which is better than nothing! How simple can we make this matrix? Let $\vv_2=\stwovec xy$, then
\begin{align*}
&A\vv_2=\twovec{-8x+4y}{-25x+12y}=2\twovec xy+\twovec{-10x+4y}{-25x+10y}=2\vv_2+(-5x+2y)\vv_1\\
&\implies [\rL_A]_\beta=\begin{pmatrix}
2&-5x+2y\\0&2
\end{pmatrix}
\end{align*}
Since $\vv_2$ cannot be parallel to $\vv_1$, the only thing we cannot have is a diagonal matrix. The next best thing is for the upper right corner be 1; for instance we could choose
\[\beta=\{\vv_1,\vv_2\}=\left\{\twovec 25,\twovec 13\right\}\implies [\rL_A]_\beta=\begin{pmatrix}
2&1\\0&2
\end{pmatrix}\]
\end{example}


\begin{defn}{}{}
A \emph{Jordan block} is a square matrix of the form
\[J=\begin{pmatrix}
\lambda&1&&\\[-2pt]
&\lambda&\ddots&\\[-2pt]
&&\ddots&1\\
&&&\lambda
\end{pmatrix}\]
where all non-indicated entries are zero. Any $1\times 1$ matrix is also a Jordan block.\smallbreak
A \emph{Jordan canonical form} is a block-diagonal matrix $\diag(J_1,\ldots,J_m)$ where each $J_k$ is a Jordan block.\smallbreak
A \emph{Jordan canonical basis} for $\rT\in\cL(V)$ is a basis $\beta$ of $V$ such that $[\rT]_\beta$ is a Jordan canonical form.
\end{defn}

%In the above example, $\beta$ is a Jordan canonical basis, and $[\rL_A]_\beta$ and Jordan canonical form.\smallbreak

If a map is diagonalizable, then any eigenbasis is Jordan canonical and the corresponding Jordan canonical form is diagonal. What about more generally? Does every non-diagonalizable map have a Jordan canonical basis? If so, how can we find such?

\goodbreak

\begin{example}{}{jordan2}
It can easily be checked that
$\beta=\{\vv_1,\vv_2,\vv_3\}=\left\{\sthreevec 101,\sthreevec 120,\sthreevec 111\right\}$
is a Jordan canonical basis for 
\[A=\begin{pmatrix}
-1&2&3\\-4&5&4\\-2&1&4
\end{pmatrix}\]
(really $\rL_A\in\cL(\R^3)$). Indeed
\[A\vv_1=2\vv_1,\quad A\vv_2=3\vv_2,\quad A\vv_3=\threevec 453=\threevec{1+3}{2+3}{0+3}=\vv_2+3\vv_3\implies [\rL_A]_\beta=\begin{pmatrix}
2&0&0\\
0&3&1\\
0&0&3
\end{pmatrix}\]
\end{example}

\boldsubsubsection{Generalized Eigenvectors}

Example \ref{ex:jordan2} was easy to check, but how would we go about finding a suitable $\beta$ if we were merely given $A$? We brute-forced this in Example \ref{ex:jordan1}, but such is not a reasonable approach in general. Eigenvectors get us some of the way:
\begin{itemize}
  \item $\vv_1$ is an eigenvector in Example \ref{ex:jordan1}, but $\vv_2$ is not. 
  \item $\vv_1$ and $\vv_2$ are eigenvectors in Example \ref{ex:jordan2}, but $\vv_3$ is not.
\end{itemize}
The practical question is how to fill out a Jordan canonical basis once we have a maximal independent set of eigenvectors. We now define the necessary objects.

\begin{defn}{}{}
Suppose $\rT\in\cL(V)$ has an eigenvalue $\lambda$. Its \emph{generalized eigenspace} is
\[K_\lambda:=\{\vx\in V:(\rT-\lambda \I)^k(\vx)=\V0\text{ for some }k\in\N\}=\bigcup_{k\in\N}\cN(\rT-\lambda\I)^k\]
A \emph{generalized eigenvector} is any non-zero $\vv\in K_\lambda$.
\end{defn}

As with eigenspaces, the generalized eigenspaces of $A\in M_n(\F)$ are those of the map $\rL_A\in\cL(\F^n)$.\smallbreak

\begin{tcolorbox}[exstyle]
It is easy to check that our earlier Jordan canonical bases consist of generalized eigenvectors.
\begin{description}
  \item[\normalfont{Example \ref{ex:jordan1}}:] We have one eigenvalue $\lambda=2$. Since $(A-2I)^2=\begin{smatrix}
  0&0\\0&0
  \end{smatrix}$ is the zero matrix, every non-zero vector is a generalized eigenvector; plainly $K_2=\R^2$.

  \item[\normalfont Example \ref{ex:jordan2}:] We see that
  \[(A-2I)\vv_1=\V0,\qquad (A-3I)\vv_2=\V0,\qquad (A-3I)^2\vv_3=(A-3I)\vv_2=\V0\]
  whence $\beta$ is a basis of generalized eigenvectors. Indeed
  \[K_3=\Span\{\vv_1,\vv_2\},\quad K_2=E_2=\Span\{\vv_3\}\]
  though verifying this with current technology is a little awkward\ldots
\end{description}
\end{tcolorbox}

\goodbreak

In order to easily compute generalized eigenspaces, it is useful to invoke the main result of this section. We postpone the proof for a while due to its meatiness.

\begin{thm}{}{generalized1}
Suppose that the characteristic polynomial of $\rT\in\cL(V)$ splits over $\F$:
\[p(t)=(\lambda_1-t)^{m_1}\cdots (\lambda_k-t)^{m_k}\]
where the $\lambda_j$ are the \emph{distinct} eigenvalues of $\rT$ with algebraic multiplicities $m_j$. Then:
\begin{enumerate}\itemsep0pt
  \item For each eigenvalue; \ (a) \ $K_{\lambda}=\cN(\rT-\lambda \I)^{m}$ and \ (b) \ $\dim K_{\lambda}=m$.
	\item $V=K_{\lambda_1}\oplus\cdots\oplus K_{\lambda_k}$: there exists a basis of generalized eigenvectors.
\end{enumerate}
\end{thm}

Compare this with the statement on diagonalizability from the start of the course.\smallbreak

With regard to part 2; we shall eventually be able to \emph{choose} this to be a Jordan canonical basis. In conclusion: a map has a Jordan canonical basis if and only if its characteristic polynomial splits.



\begin{examples}{}{}
\exstart Observe how Example \ref{ex:jordan2} works in this language:
\begin{gather*}
A=\begin{pmatrix}
-1&2&3\\-4&5&4\\-2&1&4
\end{pmatrix}\implies p(t)=(2-t)^{\textcolor{blue}{1}}(3-t)^{\textcolor{red}{2}}\\
K_2=\cN(A-2I)^{\textcolor{blue}{1}}=\Span\threevec 101\implies\dim K_2={\textcolor{blue}{1}}\\
K_3=\cN(A-3I)^{\textcolor{red}{2}}=\cN\begin{pmatrix}
2&-1&-1\\0&0&0\\2&-1&-1
\end{pmatrix} =\Span\left\{\threevec 120,\threevec 111\right\}\implies\dim K_3={\textcolor{red}{2}}\\
\R^3=\textcolor{blue}{K_2}\oplus \textcolor{red}{K_3}
\end{gather*}
\begin{enumerate}\setcounter{enumi}{1}
  \item We find the generalized eigenspaces of the matrix $A=\begin{smatrix}
5&2&-1\\0&0&0\\9&6&-1
\end{smatrix}$\\
The characteristic polynomial is
\[p(t)=\det(A-\lambda I)=-t\begin{vmatrix}
5-t&-1\\9&-1-t
\end{vmatrix} =-t(t^2-5t+t-5+9)=-(0-t)^{\textcolor{blue}{1}}(2-t)^{\textcolor{red}{2}}\]
\begin{itemize}
  \item $\textcolor{blue}{\lambda=0}$ has multiplicity \textcolor{blue}{1}; indeed $K_0=\cN(A-0I)^{\textcolor{blue}{1}}=\cN(A)=\Span\sthreevec 1{-1}3$ is just the eigenspace $E_0$.
  \item $\textcolor{red}{\lambda=2}$ has multiplicity $\textcolor{red}{2}$,
  \[K_2=\cN(A-2I)^{\textcolor{red}{2}} =\cN\begin{pmatrix}
  3&2&-1\\0&-2&0\\9&6&-3
  \end{pmatrix}^2 =\cN\begin{pmatrix}
  0&-4&0\\0&4&0\\0&-12&0
  \end{pmatrix}=\Span\left\{\threevec 100,\threevec 001\right\}\]
  In this case the corresponding eigenspace is \emph{one}-dimensional, $E_2=\Span\sthreevec 103\lneq K_2$, and the matrix is non-diagonalizable.
\end{itemize}
Observe also that $\R^3=\textcolor{blue}{K_0}\oplus \textcolor{red}{K_2}$ in accordance with the Theorem.
\end{enumerate}
\end{examples}

\subsubsection*{Properties of Generalized Eigenspaces and the Proof of Theorem \ref{thm:generalized1}}

A lot of work is required to justify our main result. Feel free to skip the proofs at first reading.

\begin{lemm}{}{gevectorbasic}
Let $\lambda$ be an eigenvalue of $\rT\in\cL(V)$. Then:
\begin{enumerate}
  \item $E_\lambda$ is a subspace of $K_\lambda$, which is itself a subspace of $V$.
  \item $K_\lambda$ is $\rT$-invariant.
  \item Suppose $K_\lambda$ is finite-dimensional and $\mu\neq \lambda$. Then:
  \begin{enumerate}
    \item $K_\lambda$ is $(\rT-\mu \I)$-invariant and the restriction of $\rT-\mu\I$ to $K_\lambda$ is an isomorphism.
    \item If $\mu$ is another eigenvalue, then $K_\lambda\cap K_\mu=\{\V0\}$. In particular $K_\lambda$ contains no eigenvectors other than those in $E_\lambda$. 
	\end{enumerate}
\end{enumerate}
\end{lemm}

\begin{proof}
\begin{enumerate}
  \item These are an easy exercise.
  
  \item Let $\vx\in K_\lambda$, then $\exists k$ such that $(\rT-\lambda \I)^k(\vx)=\V0$. But then
  \begin{align*}
  (\rT-\lambda \I)^k\bigl(\rT(\vx)\bigr)&=(\rT-\lambda \I)^{k}\bigl(\rT(\vx)-\lambda \vx+\lambda\vx\bigr)\\
  &=(\rT-\lambda \I)^{k+1}(\vx)+\lambda(\rT-\lambda \I)^k(\vx) =\V0
  \end{align*}
  Otherwise said, $\rT(\vx)\in K_\lambda$.
  
  \item\begin{enumerate}
    \item Let $\vx\in K_\lambda$. Part 2 tells us that
    \[(\rT-\mu\I)(\vx)=\rT(\vx)-\mu\vx\in K_\lambda\]
    whence $K_\lambda$ is $(\rT-\mu \I)$-invariant.\smallbreak
    Suppose, for a contradiction, that $\rT-\mu\I$ is \emph{not injective} on $K_\lambda$. Then
		\[\exists\vy\in K_\lambda\setminus\{\V0\}\text{ such that }(\rT-\mu \I)(\vy)=\V0\]
		Let $k\in\N$ be minimal such that $(\rT-\lambda \I)^k(\vy)=\V0$ and let $\vz=(\rT-\lambda \I)^{k-1}(\vy)$. Plainly $\vz\neq\V0$, for otherwise $k$ is not minimal. Moreover,
		\[(\rT-\lambda\I)(\vz)=(\rT-\lambda \I)^k(\vy)=\V0\implies \vz\in E_\lambda\]
		Since $\rT-\mu\I$ and $\rT-\lambda\I$ commute, we can also compute the effect of $\rT-\mu\I$;
		\[
		(\rT-\mu \I)(\vz)=(\rT-\mu\I)(\rT-\lambda \I)^{k-1}(\vy) =(\rT-\lambda \I)^{k-1}(\rT-\mu \I)(\vx)=\V0
		\]
		which says that $\vz$ is an eigenvector in $E_\mu$; if $\mu$ isn't an eigenvalue, then we already have our contradiction! Even if $\mu$ is an eigenvalue, $E_\mu\cap E_\lambda=\{\V0\}$ provides the desired contradiction.\smallbreak
		We conclude that $(\rT-\mu\I)_{K_\lambda}\in\cL(K_\lambda)$ is injective. Since $\dim K_\lambda<\infty$, the restriction is automatically an isomorphism.
		\item	This is another exercise.\qedhere
  \end{enumerate} 
\end{enumerate}
\end{proof}

\goodbreak

Now to prove Theorem \ref{thm:generalized1}: remember that the characteristic polynomial of $\rT$ is assumed to split.



\begin{proof}
\begin{description}
\item[\normalfont{(Part 1(a))}] Fix an eigenvalue $\lambda$. By definition, we have $\cN(\rT-\lambda \I)^m\le K_\lambda$.\smallbreak
	For the converse, parts 2 and 3 of the Lemma tell us (why?) that
	\[p_\lambda(t)=(\lambda-t)^{\dim K_\lambda}\quad\text{from which}\quad \dim K_{\lambda}\le m \tag{$\ast$}\]
	By the Cayley--Hamilton Theorem, $\at\rT{K_\lambda}$ satisfies its characteristic polynomial, whence
	\[\forall \vx\in K_\lambda,\  \left(\lambda \I-\rT\right)^{\dim K_\lambda}(\vx)=\V0 \implies K_\lambda\le \cN(\rT-\lambda \I)^m\]
	%We've therefore established that $K_\lambda=\cN(\rT-\lambda \I)^m$.
	
	\item[\normalfont{(Parts 1(b) and 2)}] We prove simultaneously by induction on the number of distinct eigenvalues of $\rT$.\smallbreak
(\emph{Base case})\quad If $\rT$ has only one eigenvalue, then $p(t)=(\lambda-t)^m$. Another appeal to Cayley--Hamilton says $(\rT-\lambda \I)^m(\vx)=\V0$ for all $\vx\in V$. Thus $V=K_\lambda$ and $\dim K_\lambda=m$.\smallbreak
(\emph{Induction step})\quad Fix $k$ and suppose the results hold for maps with $k$ distinct eigenvalues. Let $\rT$ have distinct eigenvalues $\lambda_1,\ldots,\lambda_k,\mu$, with multiplicities $m_1,\ldots,m_k,m$ respectively. Define\footnotemark
\[\smash{W=\cR(\rT-\mu\I)^{m}}\]
The subspace $W$ has the following properties, the first two of which we leave as exercises:\vspace{-3pt}
\begin{itemize}\itemsep1pt
  \item $W$ is $\rT$-invariant.
  \item $W\cap K_\mu=\{\V0\}$ so that $\mu$ is not an eigenvalue of the restriction $\at\rT W$.
  \item Each $K_{\lambda_j}\le W$: since $(\rT-\mu\I)_{K_{\lambda_j}}$ is an isomorphism (Lemma part 3), we can invert,
  \[\smash{\vx\in K_{\lambda_j}\implies \vx=(\rT-\mu\I)^{m}\Bigl((\rT-\mu\I)_{K_{\lambda_j}}^{-1}\Bigr)^m(\vx) \in \cR(\rT-\mu\I)^{m}=W}\]
  We conclude that $\lambda_j$ is an eigenvalue of the restriction $\at\rT W$ with generalized eigenspace $K_{\lambda_j}$.
\end{itemize}
Since $\at\rT W$ has $k$ distinct eigenvalues, the induction hypotheses apply:
\[\smash{W=K_{\lambda_1}\oplus\cdots\oplus K_{\lambda_{k}}\quad\text{and}\quad p_{W}(t)=(\lambda_1-t)^{\dim K_{\lambda_1}}\cdots (\lambda_{k}-t)^{\dim K_{\lambda_{k}}}}\]
Since $W\cap K_\mu=\{\V0\}$ it is enough finally to use the rank--nullity theorem and count dimensions:\vspace{-8pt}
\begin{align*}
\dim V&=\Rank (\rT-\mu\I)^{m}+\Null (\rT-\mu\I)^{m}=\dim W+\dim K_\mu =\smash[b]{\sum_{j=1}^k}\dim K_{\lambda_j}+\dim K_\mu\\
&\overset{(\ast)}{\le} m_1+\cdots+m_k+m =\operatorname{deg}(p(t))=\dim V
\end{align*}
The inequality is thus an \emph{equality}; each $\dim K_{\lambda_j}=m_j$ and $\dim K_\mu=m$. We conclude that
\[V=K_{\lambda_1}\oplus\cdots\oplus K_{\lambda_{k}}\oplus K_\mu\]
which completes the induction step and thus the proof. Whew!\qedhere
\end{description}
\end{proof}

\footnotetext{This is yet another argument where we consider a suitable subspace to which we can apply an induction hypothesis; recall the spectral theorem, Schur's lemma, bilinear form diagonalization, etc. Theorem \ref{thm:jordanbasis} will provide one more!}


\goodbreak

\boldsubsubsection{Cycles of Generalized Eigenvectors}

By Theorem \ref{thm:generalized1}, for every linear map whose characteristic polynomial splits there exists generalized eigenbasis. This isn't the same as a Jordan canonical basis, but we're very close!

\begin{example}{}{}
The matrix $A=\begin{smatrix}
5&1&0\\0&5&1\\0&0&5
\end{smatrix}\in M_3(\R)$ is a single Jordan block, whence there is a single generalized eigenspace $K_5=\R^3$ and the standard basis $\epsilon=\{\ve_1,\ve_2,\ve_3\}$ is Jordan canonical.\smallbreak
The crucial observation for what follows is that one of these vectors $\ve_3$ \emph{generates} the others via repeated applications of $A-5I$:
\[\ve_2=(A-5I)\ve_3,\qquad \ve_1=(A-5I)\ve_2=(A-5I)^2\ve_3\]
\end{example}


\begin{defn}{}{}
A \emph{cycle of generalized eigenvectors} for a linear operator $\rT$ is a set
\[\beta_\vx:=\smash{\left\{(\rT-\lambda \I)^{k-1}(\vx),\ldots,(\rT-\lambda \I)(\vx),\vx\right\}}\]
where the \emph{generator} $\vx\in K_\lambda$ is non-zero and $k$ is minimal such that $(\rT-\lambda \I)^k(\vx)=\V0$.
\end{defn}

Note that the first element $(\rT-\lambda\I)^{k-1}(\vx)$ is an \emph{eigenvector}.

Our goal is to show that $K_\lambda$ has a basis consisting of cycles of generalized eigenvectors; putting these together results in a Jordan canonical basis.

\begin{lemm}{}{gencyclesbasic}
Let $\beta_\vx$ be a cycle of generalized eigenvectors of $\rT$ with length $k$. Then:
\begin{enumerate}%\itemsep2pt
  \item $\beta_\vx$ is linearly independent and thus a basis of $\Span\beta_\vx$.
  \item $\Span\beta_\vx$ is $\rT$-invariant. With respect to $\beta_\vx$, the matrix of the restriction of $\rT$ is the $k\times k$ Jordan block $[\rT_{\Span\beta_\vx}]_{\beta_\vx}=\begin{smatrix}
  \lambda&1&&\\[-6pt]
  &\lambda&\ddots&\\[-6pt]
  &&\ddots&1\\
  &&&\lambda
  \end{smatrix}$.
\end{enumerate}
\end{lemm}

In what follows, it will be useful to consider the linear map $\rU=\rT-\lambda\I$. Note the following:
\begin{itemize}%\itemsep2pt
  \item The nullspace of $\rU$ is the \emph{eigenspace}: $\cN(\rU)=E_\lambda\le K_\lambda$.
  \item $\rT$ commutes with $\rU$: that is $\rT\rU=\rU\rT$.
  \item $\beta_\vx=\{\rU^{k-1}(\vx),\ldots,\rU(\vx),\vx\}$; that is, $\Span\beta_\vx=\ip{\vx}$ is the $\rU$-cyclic subspace generated by $\vx$.
\end{itemize}

\begin{proof}
\begin{enumerate}%\itemsep2pt
  \item Feed the linear combination $\sum_{j=0}^{k-1}a_j\rU^j(\vx)=\V0$ to $\rU^{k-1}$ to obtain
  \[a_0\rU^{k-1}(\vx)=\V0\implies a_0=0\]
  Now feed the same combination to $\rU^{k-2}$, etc., to see that all coefficients $a_j=0$.
  \item Since $\rT$ and $\rU$ commute, we see that
  \[\rT\bigl(\rU^j(\vx)\bigr)=\rU^j\bigl(\rT(\vx)\bigr)=\rU^j\bigl((\rU+\lambda\I)(\vx)\bigr)=\rU^{j+1}(\vx)+\lambda\rU^j(\vx)\in\Span\beta_\vx\]
  This justifies both $\rT$-invariance and the Jordan block claim!\qedhere
\end{enumerate}
\end{proof}

\goodbreak

The basic approach to finding a Jordan canonical basis is to find the generalized eigenspaces and play with cycles until you find a basis for each $K_\lambda$. \emph{Many} choices of canonical basis exist for a given map! We'll consider a more systematic method in the next section.

\begin{examples}{}{jordansimple1}
\exstart The characteristic polynomial of $A=\begin{smatrix}
	1&0&2\\0&1&6\\6&-2&1
\end{smatrix}\in M_3(\R)$ splits:
	\[p(t)=(1-t)\smash{\begin{vmatrix}
		1-t&6\\-2&1-t
		\end{vmatrix}} +2\smash{\begin{vmatrix}
		0&1-t\\6&-2
	\end{vmatrix}} =(1-t)\left((1-t)^2+12-12\right)=(1-t)^3\]
	With only one eigenvalue we see that $K_1=\R^3$. Simply choose any vector in $\R^3$ and see what $U=A-I$ does to it! For instance, with $\vx=\ve_1$,
	\[\beta_\vx=\left\{\rU^2\sthreevec 100,\rU\sthreevec 100,\sthreevec 100\right\}=\left\{\sthreevec{12}{36}0,\sthreevec 006,\sthreevec 100\right\}\]
	provides a Jordan canonical basis of $\R^3$. We conclude
	\[A=QJQ^{-1}=
	\begin{smatrix}
		12&0&1\\36&0&0\\0&6&0
		\end{smatrix}
	\begin{smatrix}
		1&1&0\\0&1&1\\0&0&1
	\end{smatrix}
	\begin{smatrix}
		12&0&1\\36&0&0\\0&6&0
	\end{smatrix}^{-1}\] 
	In practice, almost any choice of $\vx\in\R^3$ will generate a cycle of length three!

\begin{enumerate}\setcounter{enumi}{1}
  \item The matrix $B=\begin{smatrix}
7&1&-4\\0&3&0\\8&1&-5
\end{smatrix}\in M_3(\R)$ has characteristic equation
\[p(t)=(3-t)(t^2-2t-3)=-(t+1)^{\textcolor{blue}1}(t-3)^{\textcolor{red}2}\]
$\dim K_{-1}=\textcolor{blue}1\implies K_{-1}=E_{-1}=\Span\smash{\sthreevec 102}$, spanned by a cycle of length one.\smallbreak
Since $\dim K_3=\textcolor{red}2$, we have
  \[K_3=\cN(B-3I)^{\textcolor{red}2}=\cN\smash{\begin{smatrix}
  4&1&-4\\0&0&0\\8&1&-8
  \end{smatrix}^2=
  \cN\begin{smatrix}
  -16&0&16\\0&0&0\\-32&0&32
  \end{smatrix}=\Span\left\{\sthreevec 101,\sthreevec 010\right\}}\]
  This is spanned by a cycle of length two: $\sthreevec 101$ is an eigenvector and
  \[\smash{\sthreevec 101=(B-3I)\sthreevec 010}\]
  We conclude that $\beta=\left\{\sthreevec 102,\sthreevec 101,\sthreevec 010\right\}$ is a Jordan canonical basis for $B$, and that
	\[B=QJQ^{-1}=
		\begin{smatrix}
		1&1&0\\0&0&1\\2&1&0
	\end{smatrix}
	\begin{smatrix}
		-1&0&0\\0&3&1\\0&0&3
	\end{smatrix}
	\begin{smatrix}
		1&1&0\\0&0&1\\2&1&0
	\end{smatrix}^{-1}\]
 
	\item Let $\rT=\diff x$ on $P_3(\R)$. With respect to the standard basis $\epsilon=\{1,x,x^2,x^3\}$,
	\[A=[\rT]_\epsilon =\begin{smatrix}
		0&1&0&0\\0&0&2&0\\0&0&0&3\\0&0&0&0
	\end{smatrix}\]
	With only one eigenvalue $\lambda=0$, we have a single generalized eigenspace $K_0=P_3(\R)$. It is easy to check that $f(x)=x^3$ generates a cycle of length three and thus a Jordan canonical basis:
	\[\beta=\{6,6x,3x^2,x^3\}\implies[\rT]_\beta =\begin{smatrix}
		0&1&0&0\\0&0&1&0\\0&0&0&1\\0&0&0&0
	\end{smatrix}
	=
	\begin{smatrix}
		6&0&0&0\\0&6&0&0\\0&0&3&0\\0&0&0&1
	\end{smatrix}^{-1}
	\begin{smatrix}
		0&1&0&0\\0&0&2&0\\0&0&0&3\\0&0&0&0
	\end{smatrix}
	\begin{smatrix}
		6&0&0&0\\0&6&0&0\\0&0&3&0\\0&0&0&1
	\end{smatrix}\]
\end{enumerate}
\end{examples}

Our final results state that this process works generally.

\begin{thm}{}{jordanbasis}
Let $\rT\in\cL(V)$ have an eigenvalue $\lambda$. If $\dim K_\lambda<\infty$, then there exists a basis $\beta_\lambda=\beta_{\vx_1}\cup\cdots\cup\beta_{\vx_n}$ of $K_\lambda$ consisting of finitely many linearly independent cycles.
\end{thm}

Intuition suggests that we create cycles $\beta_{\vx_j}$ by starting with a basis of the eigenspace $E_\lambda$ and extending backwards: for each $\vx$, if $\vx=(\rT-\lambda\I)(\vy)$, then $\vx\in\beta_\vy$; now repeat until you have a maximum length cycle. This is essentially what we do, though a sneaky induction is required to make sure we keep track of everything and guarantee that the result really is a basis of $K_\lambda$.

\begin{proof}
We prove by induction on $m=\dim K_\lambda$.\smallbreak
(\emph{Base case})\quad  If $m=1$, then $K_\lambda=E_\lambda=\Span\vx$ for some eigenvector $\vx$. Plainly $\{\vx\}=\beta_\vx$.\smallbreak
(\emph{Induction step})\quad  Fix $m\ge 2$. Write $n=\dim E_\lambda\le m$ and $\rU=(\rT-\lambda\I)_{K_\lambda}$. 
\begin{enumerate}\renewcommand{\labelenumi}{(\roman{enumi})}\itemsep0pt
  \item For the induction hypothesis, suppose \emph{every} generalized eigenspace with dimension $<m$ (for \emph{any} linear map!) has a basis consisting of independent cycles of generalized eigenvectors.
  \item Define $W=\cR(\rU)\cap E_\lambda$: that is
  \[\vw\in W\iff\begin{cases}
  \rU(\vw)=\V0\quad\text{and}\\
  \vw=\rU(\vv)\text{ for some }\vv\in K_\lambda
  \end{cases}\]
  Let $k=\dim W$, choose a complementary subspace $X$ such that $E_\lambda=W\oplus X$ and select a basis $\{\vx_{k+1},\ldots,\vx_n\}$ of $X$. If $k=0$, the induction step is finished (why?). Otherwise we continue\ldots
  \item The calculation in the proof of Lemma \ref{lemm:gencyclesbasic} (take $j=1$) shows that $\cR(\rU)$ is $\rT$-invariant; it is therefore the single generalized eigenspace $\tilde{K}_\lambda$ of $\rT_{\cR(\rU)}$.
  \item By the rank--nullity theorem,
	\[\dim \cR(\rU)=\Rank\rU=\dim K_\lambda-\Null\rU=m-\dim E_\lambda <m\]
	By the induction hypothesis, $\cR(\rU)$ has a basis of independent cycles. Since the last non-zero element in each cycle is an eigenvector, this basis consists of $k$ distinct cycles $\beta_{\hat\vx_1}\cup\cdots\cup\beta_{\hat\vx_k}$ whose terminal vectors form a basis of $W$.
  \item Since each $\hat\vx_j\in\cR(\rU)$, there exist vectors $\vect xk$ such that $\hat\vx_j=\rU(\vx_j)$. Including the length-one cycles generated by the basis of $X$, the cycles $\beta_{\vx_1},\ldots,\beta_{\vx_n}$ now contain
  \[\dim\cR(\rU)+k+(n-k)=\Rank\rU+\Null\rU=m\]
  vectors. We leave as an exercise the verification that these vectors are linearly independent.\qedhere
\end{enumerate} 
\end{proof}

\begin{cor}{}{jordanbasiscor}
Suppose that the characteristic polynomial of $\rT\in\cL(V)$ splits (necessarily $\dim V<\infty$). Then there exists a Jordan canonical basis, namely the union of bases $\beta_\lambda$ from Theorem \ref{thm:jordanbasis}.
\end{cor}

\begin{proof}
By Theorem \ref{thm:generalized1}, $V$ is the direct sum of generalized eigenspaces. By the previous result, each $K_\lambda$ has a basis $\beta_\lambda$ consisting of finitely many cycles. By Lemma \ref{lemm:gencyclesbasic}, the matrix of $\rT_{K_\lambda}$ has Jordan canonical form with respect to $\beta_\lambda$. It follows that $\beta=\bigcup\beta_\lambda$ is a Jordan canonical basis for $\rT$.
\end{proof}

\goodbreak

\begin{exercises}
\exstart For each matrix, find the generalized eigenspaces $K_\lambda$, find bases consisting of unions of disjoint cycles of generalized eigenvectors, and thus find a Jordan canonical form $J$ and invertible $Q$ so that the matrix may be expressed as $QJQ^{-1}$.
\begin{enumerate}\setcounter{enumi}{1}
  \item[]\begin{enumerate}
    \item $A=\begin{pmatrix}
    1&1\\-1&3
    \end{pmatrix}$
    \qquad
    (b)\ \ $B=\begin{pmatrix}
    1&2\\3&2
    \end{pmatrix}$
		\qquad
    (c)\ \ $C=\begin{pmatrix}
    11&-4&-5\\21&-8&-11\\3&-1&0
    \end{pmatrix}$
    \item[(d)] $D=\begin{pmatrix}
    2&1&0&0\\0&2&1&0\\0&0&3&0\\0&1&-1&3
    \end{pmatrix}$
	\end{enumerate}
	
	\item If $\beta=\{\vect vn\}$ is a Jordan canonical basis, what can you say about $\vv_1$? Briefly explain why the linear map $\rL_A\in\cL(\R^2)$ where $A=\begin{smatrix}
	0&-1\\1&0
	\end{smatrix}$ has no Jordan canonical form.
	
	\item Find a Jordan canonical basis for each linear map $\rT$:
	\begin{enumerate}
	  \item $\rT\in\cL(P_2(\R))$ defined by $\rT(f(x))=2f(x)-f'(x)$
	  \item $\rT(f)=f'$ defined on $\Span\{1,t,t^2,e^t,te^t\}$
	  \item $\rT(A)=2A+A^T$ defined on $M_2(\R)$
	\end{enumerate}
	
	\item In Example \hyperref[ex:jordansimple1]{\ref*{ex:jordansimple1}.1}, suppose $\vx=\sthreevec abc$. Show that almost any choice of $a,b,c$ produces a Jordan canonical basis $\beta_\vx$.

  
  \item We complete the proof of Lemma \ref{lemm:gevectorbasic}.
  \begin{enumerate}
    \item Prove part 1: that $E_\lambda\le K_\lambda\le V$.
    \item Verify that $\rT-\mu\I$ and $\rT-\lambda\I$ commute.
    \item Prove part 3(b): generalized eigenspaces for distinct eigenvalues have trivial intersection.
  \end{enumerate}
  
  \item Consider the induction step in the proof of Theorem \ref{thm:generalized1}.
  \begin{enumerate}
    \item Prove that $W$ is $\rT$-invariant.
    \item Explain why $W\cap K_\mu=\{\V0\}$.
    \item The assumption $p_W(t)=(\lambda_1-t)^{\dim K_{\lambda_1}}\cdots (\lambda_k-t)^{\dim K_{\lambda_k}}$ near the end of the proof is the induction hypothesis for part 1(b). Why can't we also assume that $\dim K_{\lambda_j}=m_j$ and thus tidy the inequality argument near the end of the proof?
  \end{enumerate}
  
  \item We finish some of the details of Theorem \ref{thm:jordanbasis}.
  \begin{enumerate}
    \item In step (ii), suppose $\dim W=k=0$. Explain why $\{\vx_1,\ldots,\vx_n\}$ is in fact a basis of $K_\lambda$, so that the rest of the proof is unnecessary.
    \item In step (v), prove that the $m$ vectors in the cycles $\beta_{\vx_1},\ldots,\beta_{\vx_n}$ are linearly independent.\smallbreak
    (\emph{Hint: model your argument on part 1 of Lemma \ref{lemm:gencyclesbasic}})
  \end{enumerate}
 
\end{enumerate}
\end{exercises}


\clearpage


\subsection{Cycle Patterns and the Dot Diagram}


In this section we obtain a useful result that helps us compute Jordan forms more efficiently and systematically. To give us some clues how to proceed, here is a lengthy example.

\begin{example}{}{dotdiagmotiv}
Precisely three Jordan canonical forms $A,B,C\in M_3(\R)$ correspond to the characteristic polynomial $p(t)=(5-t)^3$:
\[A=\begin{pmatrix}
5&0&0\\
0&5&0\\
0&0&5
\end{pmatrix}
\qquad
B=\begin{pmatrix}
5&1&0\\
0&5&0\\
0&0&5
\end{pmatrix}
\qquad
C=\begin{pmatrix}
5&1&0\\
0&5&1\\
0&0&5
\end{pmatrix}\]
In all three cases the standard basis $\beta=\{\ve_1,\ve_2,\ve_3\}$ is Jordan canonical, so how do we distinguish things? By considering the \emph{\textcolor{blue}{number}} and \emph{\textcolor{red}{lengths}} of the cycles of generalized eigenvectors.
\begin{itemize}
  \item $A$ has eigenspace $E_5=K_5=\R^3$. Since $(A-5I)\vv=\V0$ for all $\vv\in\R^3$, we have maximum \textcolor{red}{cycle-length one}. We therefore need \textcolor{blue}{three distinct cycles} to construct a Jordan basis, e.g.
  \[\beta_{\ve_1}=\{\ve_1\},\quad \beta_{\ve_2}=\{\ve_2\},\quad \beta_{\ve_3}=\{\ve_3\} \implies \beta=\beta_{\ve_1}\cup\beta_{\ve_2}\cup\beta_{\ve_3}=\{\ve_1,\ve_2,\ve_3\}\]
  
  \item $B$ has eigenspace $E_5=\Span\{\ve_1,\ve_3\}$. By computing
  \[\vv=\threevec abc\implies (B-5I)\vv=\threevec{b}00 \implies (B-5I)^2\vv=\V0\]
  we see that $\beta_\vv$ is a cycle with \textcolor{red}{maximum length two}, provided $b\neq 0$ ($\vv\not\in E_5$). We therefore need \textcolor{blue}{two distinct cycles,} of lengths \textcolor{red}{two} and \textcolor{red}{one}, to construct a Jordan basis, e.g.
  \[\beta_{\ve_2}=\bigl\{(B-5I)\ve_2,\,\ve_2\bigr\}=\{\ve_1,\ve_2\},\quad \beta_{\ve_3}=\{\ve_3\}\implies \beta=\beta_{\ve_2}\cup\beta_{\ve_3}=\{\ve_1,\ve_2,\ve_3\}\]
  
  \item $C$ has eigenspace $E_5=\Span\ve_1$. This time
  \[\vv=\threevec abc\implies (C-5I)\vv=\threevec{b}c0,\quad (C-5I)^2\vv=\threevec c00,\quad (C-5I)^3\vv=\V0\]
  generates a cycle with \textcolor{red}{maximum length two} provided $c\neq 0$. Indeed this cycle is a Jordan basis, so \textcolor{blue}{one cycle} is all we need:
  \[\beta=\beta_{\ve_3}=\bigl\{(C-5I)^2\ve_3,\,(C-5I)\ve_3,\,\ve_3\bigr\}=\{\ve_1,\ve_2,\ve_3\}\]
\end{itemize}
\end{example}

Why is the example relevant? Suppose that $\dim_\R V=3$ and that $\rT\in\cL(V)$ has characteristic polynomial $p(t)=(5-t)^3$. Theorem \ref{thm:jordanbasis} tells us that $\rT$ has a Jordan canonical form, and that is is moreover one of the above matrices $A,B,C$. Our goal is to develop a method whereby the pattern of cycle-lengths can be determined, thus allowing us to be able to discern which Jordan form is correct. As a side-effect, this will also demonstrate that the pattern of cycle lengths for a given $\rT$ is independent of the Jordan basis so that, up to some reasonable restriction, the Jordan \emph{form} of $\rT$ is unique. To aid us in this endeavor, we require some terminology\ldots

\begin{defn}{}{dotdiag}
Let $V$ be finite dimensional and $K_\lambda$ a generalized eigenspace of $\rT\in\cL(V)$. Following the Theorem \ref{thm:jordanbasis}, assume that $\beta_\lambda=\beta_{\vx_1}\cup\cdots\cup\beta_{\vx_n}$ is a Jordan canonical basis of $\at\rT{K_\lambda}$, where the cycles are arranged in \emph{non-increasing} length. That is:
\begin{enumerate}%\itemsep2pt
  \item $\beta_{\vx_j}=\{(\rT-\lambda \I)^{k_j-1}(\vx_j),\ldots,\vx_j\}$ has length $k_j$, and
  \item $k_1\ge k_2\ge\cdots\ge k_n$
\end{enumerate}
The \emph{dot diagram} of $\at\rT{K_\lambda}$ is a representation of the elements of $\beta_\lambda$, one dot for each vector: the $j\th$ column represents the elements of $\beta_{\vx_j}$ arranged vertically with $\vx_j$ at the bottom.
\end{defn}

Given a linear map, our eventual goal is to identify the dot diagram as an intermediate step in the computation of a Jordan basis. First, however, we observe how the conversion of dot diagrams to a Jordan form is essentially trivial.

\begin{example}{}{}
Suppose $\dim V=14$ and that $\rT\in\cL(V)$ has the following eigenvalues and dot diagrams:
  \[\begin{array}{c@{\qquad}c@{\qquad}c}
  \lambda_1=-4&\lambda_2=7&\lambda_3=12\\
  \begin{array}{cccc}
  \bullet&\bullet&\bullet&\bullet\\
  \bullet&\bullet&\\
  &&
  \end{array}
  &
  \begin{array}{cc}
  \bullet&\bullet\\
  \bullet&\bullet\\
  \bullet&
  \end{array}
  &
  \begin{array}{ccc}
  \bullet&\bullet&\bullet\\
  &&\\
  &&
  \end{array}
  \end{array}\]
  Then generalized eigenspaces of $\rT$ satisfy:
  \begin{itemize}
    \item $K_{-4}=\cN(\rT+4\I)^2$ and $\dim K_{-4}=6$;
    \item $K_{7}=\cN(\rT-7\I)^3$ and $\dim K_{7}=5$;
    \item $K_{12}=\cN(\rT-12\I)=E_{12}$ and $\dim K_{12}=3$;
  \end{itemize}
	$\rT$ has a Jordan canonical basis $\beta$ with respect to which its Jordan canonical form is
	\[[\rT]_\beta =\scalebox{0.9}{$\left(\begin{array}{@{}c@{}c@{}c@{}c@{}c@{}c@{}c@{}c@{}c@{}}
		\begin{array}{@{}cc|@{}}
			-4 & 1\\
			0 & -4\\\hline
		\end{array}&&&&&&&&\\[-0.4pt]
		&\begin{array}{|cc|}\hline
			-4 & 1\\
			0 & -4\\\hline
		\end{array}&&&&&&&\\[-0.4pt]
		&&\begin{array}{|c|}\hline
			-4\\\hline
		\end{array}&&&&&&\\[-0.4pt]
		&&&\begin{array}{|c|}\hline
			-4\\\hline
		\end{array}&&&&&\\[-0.4pt]
		&&&&\begin{array}{|ccc|}\hline
			7 & 1 & 0\\
			0 & 7 & 1\\
			0 & 0 & 7\\\hline
		\end{array}&&&&\\[-0.4pt]
		&&&&&\begin{array}{|cc|}\hline
			7 & 1\\
			0 & 7\\\hline
		\end{array}&&&\\[-0.4pt]
		&&&&&&\begin{array}{|c|}\hline
			12\\\hline
		\end{array}&&\\[-0.4pt]
		&&&&&&&\begin{array}{|c|}\hline
			12\\\hline
		\end{array}&\\[-0.4pt]
		&&&&&&&&\begin{array}{|c}\hline
			12
		\end{array}
		\end{array}\right)$}\]
	Note how the sizes of the Jordan blocks are \emph{non-increasing} within each eigenvalue. For instance, for $\lambda_1=-4$, the sequence of cycle lengths ($k_j$) is $2\ge 2\ge 1\ge 1$.
\end{example}

\goodbreak

\begin{thm}{}{}
Suppose $\beta_\lambda$ is a Jordan canonical basis of $\at\rT{K_\lambda}$ as described in Definition \ref{defn:dotdiag}, and suppose the $i^\text{th}$ row of the dot diagram has $r_i$ entries. Then:
\begin{enumerate}%\itemsep2pt
  \item For each $r\in\N$, the vectors associated to the dots in the first $r$ rows form a basis of $\cN(\rT-\lambda \I)^r$.
  \item $r_1=\Null(\rT-\lambda\I)=\dim V-\Rank(\rT-\lambda \I)$
  \item When $i>1$, $r_i=\Null(\rT-\lambda \I)^{i}-\Null(\rT-\lambda \I)^{i-1}= \Rank(\rT-\lambda \I)^{i-1}-\Rank(\rT-\lambda \I)^{i}$
\end{enumerate}
\end{thm}

\begin{example*}{\ref{ex:dotdiagmotiv} cont}{}
We describe the dot diagrams of the three matrices $A$, $B$, $C$, along with the corresponding vectors in the Jordan canonical basis $\beta$ and the values $r_i$.
\begin{itemize}
	\item[$A:$] $\begin{array}{lll}
		\bullet&\bullet&\bullet
		\end{array}
		\qquad
		\begin{array}{lll}
		\vx_1&\vx_2&\vx_3
		\end{array}
		\qquad
		\begin{array}{lll}
		\ve_1&\ve_2&\ve_3
		\end{array}$
		\begin{itemize}
		  \item[]\hspace*{-20pt}Since $A-5I$ is the zero matrix, $r_1=3-\Rank(A-5I)=3$. The dot diagram has one row, corresponding to three independent cycles of length one: $\beta=\beta_{\ve_1}\cup\beta_{\ve_2}\cup\beta_{\ve_3}$.
		\end{itemize}
		
	\item[$B:$] $\begin{array}[t]{ll}
		\bullet&\bullet\\
		\bullet&
		\end{array}
		\qquad
		\begin{array}[t]{cl}
		(B-5I)\vx_1&\vx_2\\
		\vx_1&
		\end{array}
		\qquad
		\begin{array}[t]{ll}
		\ve_1&\ve_3\\
		\ve_2
		\end{array}$
		\begin{itemize}
		  \item[]\hspace*{-20pt}Row 1:\quad $B-5I=\smash[t]{\begin{smatrix}
			0&1&0\\0&0&0\\0&0&0
		  \end{smatrix}}\implies \Rank(B-5I)=1$ and $r_1=3-1=2$. The first row $\{\ve_1,\ve_3\}$ is a basis of $E_5=\cN(B-5I)$.
			\item[]\hspace*{-20pt}Row 2:\quad $(B-5I)^2$ is the zero matrix, whence $r_2=\Rank(B-5I)-\Rank(B-5I)^2=1-0=1$.
		\end{itemize}
		The dot diagram corresponds to $\beta=\beta_{\ve_2}\cup\beta_{\ve_3}=\{\ve_1,\ve_2\}\cup\{\ve_3\}$.
		
		
	\item[$C:$] $\begin{array}[t]{lll}
		\bullet\\\bullet\\\bullet
		\end{array}
		\qquad
		\begin{array}[t]{c}
		(C-5I)^2\vx_1\\
		(C-5I)\vx_1\\
		\vx_1
		\end{array}
		\qquad
		\begin{array}[t]{l}
		\ve_1\\\ve_2\\\ve_3
		\end{array}$
		\begin{itemize}
		  \item[]\hspace*{-20pt}Row 1:\quad $C-5I=\smash[t]{\begin{smatrix}
			0&1&0\\0&0&1\\0&0&0
		  \end{smatrix}}			\implies r_1=3-\Rank(C-5I)=1$. The first row $\{\ve_1\}$ is a basis of $E_5=\cN(C-5I)$.
			\item[]\hspace*{-20pt}Row 2:\quad $(C-5I)^2=\begin{smatrix}
			0&0&1\\0&0&0\\0&0&0
			\end{smatrix}\implies r_2=\Rank(C-5I)-\Rank(C-5I)^2=2-1=1$. The first two rows $\{\ve_1,\ve_2\}$ form a basis of $\cN(C-5I)^2$.
			\item[]\hspace*{-20pt}Row 3:\quad $(C-5I)^3$ is the zero matrix, whence $r_3=\Rank(C-5I)^2-\Rank(C-5I)^3=1-0=1$.
		\end{itemize}
\end{itemize}
\end{example*}

\begin{proof}
As previously, let $\rU=\rT-\lambda\I$.
\begin{enumerate}
  \item Since each dot represents a basis vector $\rU^p(\vv_j)$, any $\vv\in K_\lambda$ may be written uniquely as a linear combination of the dots. Applying $\rU$ simply moves all the dots up a row and all dots in the top row to $\V0$. It follows that $\vv\in \cN(\rU^r)\iff$ it lies in the span of the first $r$ rows. Since the dots are linearly independent, they form a basis.
  \item By part 1, $r_1=\dim \cN(\rU)=\Null(\rT-\lambda \I)=\dim V-\Rank(\rT-\lambda \I)$.
  \item More generally,
  \begin{align*}
  r_i&=(r_1+\cdots+r_i)-(r_1+\cdots +r_{i-1})=\dim\cN(\rU^i)-\dim \cN(\rU^{i-1})\\
  &=\Null(\rU^i)-\Null(\rU^{i-1})=\Rank(\rT-\lambda \I)^{i-1}-\Rank(\rT-\lambda \I)^{i}\tag*{\qedhere}
  \end{align*}
\end{enumerate}
\end{proof}

\goodbreak

Since the ranks of maps $(\rT-\lambda\I)^i$ are independent of basis, so also is the dot diagram\ldots

\begin{cor}{}{jordanunique}
For any eigenvalue $\lambda$, the dot diagram is uniquely determined by $\rT$ and $\lambda$. If we list Jordan blocks for each eigenspace in non-increasing order, then the Jordan form of a linear map is unique up to the order of the eigenvalues.
\end{cor}

We now have a slightly more systematic method for finding Jordan canonical bases.

\begin{example}{}{}
The matrix $A=\begin{smatrix}
6&2&-4&-6\\
0&3&0&0\\
0&0&3&0\\
2&1&-2&-1
\end{smatrix}$ has characteristic equation
\[p(t)=(3-t)^2\begin{vmatrix}
6-t&-6\\2&-1-t
\end{vmatrix}
=(2-t)(3-t)^3\]
We have two generalized eigenspaces:
\begin{itemize}
  \item $K_2=E_2=\cN(A-2I)=\cN\begin{smatrix}
		4&2&-4&-6\\
		0&1&0&0\\
		0&0&1&0\\
		2&1&-2&-3
	\end{smatrix}=\Span\sfourvec 3002$. The trivial dot diagram $\bullet$ corresponds to this single eigenvector.
  \item $K_3=\cN(A-3I)^3$. To find the dot diagram, compute powers of $A-3I$:
		\begin{itemize}
		  \item[]\hspace*{-20pt}Row 1:\quad $A-3I=\begin{smatrix}
				3&2&-4&-6\\
				0&0&0&0\\
				0&0&0&0\\
				2&1&-2&-4
		  \end{smatrix}$ has rank 2 and the first row has $r_1=4-2=2$ entries.
			\item[]\hspace*{-20pt}Row 2:\quad $(A-3I)^2=\begin{smatrix}
				-3&0&0&6\\
				0&0&0&0\\
				0&0&0&0\\
				-2&0&0&4
			\end{smatrix}$ has rank 1 and the second row has $r_2=2-1=1$ entry.
		\end{itemize}
		Since we now have three dots (equalling $\dim K_3$), the algorithm terminates and the dot diagram for $K_3$ is $\begin{array}[t]{cc}
\bullet&\bullet\\\bullet&
\end{array}$\\
For the single dot in the second row, we choose something in $\cN(A-3I)^2$ which isn't an eigenvector; perhaps the simplest choice is $\vx_1=\ve_2$, which yields the two-cycle
\[\beta_{\vx_1}=\left\{(A-3I)\vx_1,\vx_1\right\} =\left\{\sfourvec 2001,\sfourvec 0100\right\}\]
To complete the first row, choose any eigenvector to complete the span: for instance $\vx_2=\sfourvec 0210$.
\end{itemize}

We now have suitable cycles and a Jordan canonical basis/form: 
\[\beta=\left\{\sfourvec 3002, \sfourvec 2001, \sfourvec 0100, \sfourvec 0210\right\},
\qquad
A=QJQ^{-1}=\begin{smatrix}
3&2&0&0\\
0&0&1&2\\
0&0&0&1\\
2&1&0&0
\end{smatrix}
\begin{smatrix}
2&0&0&0\\
0&3&1&0\\
0&0&3&0\\
0&0&0&3
\end{smatrix}
\begin{smatrix}
3&2&0&0\\
0&0&1&2\\
0&0&0&1\\
2&1&0&0
\end{smatrix}^{-1}\]
Other choices are available! For instance, if we'd chosen the two-cycle generated by $\vx_1=\ve_3$, we'd obtain a different Jordan basis but the same canonical form $J$:
\[\tilde\beta=\left\{\sfourvec 3002, \sfourvec{-4}00{-2}, \sfourvec 0010, \sfourvec 0210\right\}, \qquad
A=\begin{smatrix}
3&-4&0&0\\
0&0&0&2\\
0&0&1&1\\
2&-2&0&0
\end{smatrix}
\begin{smatrix}
2&0&0&0\\
0&3&1&0\\
0&0&3&0\\
0&0&0&3
\end{smatrix}
\begin{smatrix}
3&-4&0&0\\
0&0&0&2\\
0&0&1&1\\
2&-2&0&0
\end{smatrix}^{-1}\]  
\end{example}

\goodbreak

We do one final example for a non-matrix map.

\begin{example}{}{}
Let $\epsilon=\{1,x,y,x^2,y^2,xy\}$ and define $\rT\bigl(f(x,y)\bigr)=2\partials[f]{x}-\partials[f]{y}$ as a linear operator on $V=\Span_\R\epsilon$. The matrix and characteristic polynomial of $\rT$ is easy to compute:
\[[\rT]_\epsilon=\begin{smatrix}
0 & 2 & -1 & 0 & 0 & 0\\
0 & 0 & 0 & 4 & 0 & -1\\
0 & 0 & 0 & 0 & -2 & 2\\
0 & 0 & 0 & 0 & 0 & 0\\
0 & 0 & 0 & 0 & 0 & 0\\
0 & 0 & 0 & 0 & 0 & 0
\end{smatrix} \implies p(t)=t^6,\qquad [\rT^2]_\epsilon=\begin{smatrix}
0 & 0 & 0 & 8 & 2 & -4\\
0 & 0 & 0 & 0 & 0 & 0\\
0 & 0 & 0 & 0 & 0 & 0\\
0 & 0 & 0 & 0 & 0 & 0\\
0 & 0 & 0 & 0 & 0 & 0\\
0 & 0 & 0 & 0 & 0 & 0
\end{smatrix},\qquad [\rT^3]_\epsilon=O\]
There is only one eigenvalue $\lambda=0$ and therefore one generalized eigenspace $K_0=V$. We could keep working with matrices, but it is easy to translate the nullspaces of the matrices back to subspaces of $V$, from which the necessary data can be read off:
\begin{align*}
&\cN(\rT)=\Span\{1,x+2y,x^2+4y^2+4xy\}\qquad &&\Null \rT=3,\ \Rank \rT=3,\ r_1=3\\[2pt]
&\cN(\rT^2)=\Span\{1,x,y,x^2+2xy,2y^2+xy\}\qquad &&\Null \rT^2=5,\ \Rank \rT^2=1,\ r_2=3-1=2
\end{align*}
We now have five dots; since $\dim K_0=6$, the last row has one, and the dot diagram is\ $\begin{array}[t]{ccc}
\bullet&\bullet&\bullet\\
\bullet&\bullet&\\
\bullet&&
\end{array}$\\
Since the first two rows span $\cN(\rT^2)$, we may choose any $f_1\not\in\cN(\rT^2)$ for the final dot: $f_1=xy$ is suitable, from which the first column of the dot diagram becomes
\[\begin{matrix}
\rT^2(xy)&\bullet&\bullet\\[2pt]
\rT(xy)&\bullet&\\[2pt]
xy&&
\end{matrix}\qquad\quad\begin{matrix}
-4&\bullet&\bullet\\[2pt]
2y-x&\bullet&\\[2pt]
xy&&
\end{matrix}\]
Now choose the second dot on the second row to be anything in $\cN(\rT^2)$ such that the first two rows span $\cN(\rT^2)$: this time $f_2=x^2-4y^2$ is suitable, and the diagram becomes:
\[\begin{matrix}
\rT^2(xy)&\rT(x^2-4y^2)&\bullet\\[2pt]
\rT(xy)&x^2-4y^2&\\[2pt]
xy&&
\end{matrix}\qquad\quad\begin{matrix}
-4&4x+8y&\bullet\\[2pt]
2y-x&x^2-4y^2&\\[2pt]
xy&&
\end{matrix}\]
The final dot is now chosen so that the first row spans $\cN(\rT)$: this time $f_3=x^2+4y^2+4xy$ works. The result is a Jordan canonical basis and form for $\rT$
\[\beta=\bigl\{-4,\ 2y-x,\ xy,\ 4x+8y,\ x^2-4y^2,\ x^2+4y^2+4xy\bigr\},\quad\qquad J=[\rT]_\beta=
% \begin{smatrix}
% 0 & 1 & 0 & 0 & 0 & 0\\
% 0 & 0 & 1 & 0 & 0 & 0\\
% 0 & 0 & 0 & 0 & 0 & 0\\
% 0 & 0 & 0 & 0 & 1 & 0\\
% 0 & 0 & 0 & 0 & 0 & 0\\
% 0 & 0 & 0 & 0 & 0 & 0
% \end{smatrix}
% \left(\scriptsize\begin{array}{@{}c@{}c@{}c@{}}
% 		\begin{array}{@{\ }c@{\,}c@{\,}c@{\ }|}
% 			0 & 1 & 0\\
% 			0 & 0 & 1\\
% 			0 & 0 & 0\\\hline
% 		\end{array}&&\\
% 		&\begin{array}{|@{\ }c@{\,}c@{\ }|}\hline
% 			0 & 1\\
% 			0 & 0\\\hline
% 		\end{array}&\\
% 		&&\begin{array}{|@{\ }c@{\ }}\hline
% 			0
% 		\end{array}
% 	\end{array}\right)
	\scalebox{0.8}{$\left(\begin{array}{@{}c@{}c@{}c@{}}
		\begin{array}{@{\ }ccc@{\ }|}
			0 & 1 & 0\\
			0 & 0 & 1\\
			0 & 0 & 0\\\hline
		\end{array}&&\\[-0.4pt]
		&\begin{array}{|@{\ }cc@{\ }|}\hline
			0 & 1\\
			0 & 0\\\hline
		\end{array}&\\[-0.4pt]
		&&\begin{array}{|@{\ }c@{\ }}\hline
			0
		\end{array}
		\end{array}\right)$}
\]
As previously, many other choices of cycle-generators $f_1,f_2,f_3$ are available; while these result in different Jordan canonical bases, Corollary \ref{cor:jordanunique} assures us that we'll always obtain the same canonical form $J$.
\end{example}

\goodbreak

\begin{exercises}
\exstart Let $\rT$ be a linear operator whose characteristic polynomial splits. Suppose the eigenvalues and the dot diagrams for the generalized eigenspaces $K_{\lambda_i}$ are as follows:
  \[\begin{array}{c@{\qquad}c@{\qquad}c}
  \lambda_1=2&\lambda_2=4&\lambda_3=-3\\
  \begin{array}{ccc}
  \bullet&\bullet&\bullet\\
  \bullet&\bullet&\\
  \bullet&&
  \end{array}
  &
  \begin{array}{cc}
  \bullet&\bullet\\
  \bullet&\\
  \bullet&
  \end{array}
  &
  \begin{array}{cc}
  \bullet&\bullet\\
  &\\
  &
  \end{array}
  \end{array}\]
  Find the Jordan form $J$ of $\rT$.
  
\begin{enumerate}\setcounter{enumi}{1}
	\item Suppose $\rT$ has Jordan canonical form
	\[J=\scalebox{0.8}{$\left(\begin{array}{@{}c@{}c@{}c@{}c@{}}
		\begin{array}{@{\ }ccc@{\ }|}
			2 & 1 & 0\\
			0 & 2 & 1\\
			0 & 0 & 2\\\hline
		\end{array}&&&\\[-0.4pt]
		&\begin{array}{|@{\ }cc@{\ }|}\hline
			2 & 1\\
			0 & 2\\\hline
		\end{array}&&\\[-0.4pt]
		&&\begin{array}{|@{\ }c@{\ }|}\hline
			3\\\hline
		\end{array}&\\[-0.4pt]
		&&&\begin{array}{|@{\ }c@{\ }}\hline
			3
		\end{array}
		\end{array}\right)$}
% 	\begin{smatrix}
% 	2 & 1 & 0 & 0 & 0 & 0 & 0\\
% 	0 & 2 & 1 & 0 & 0 & 0 & 0\\
% 	0 & 0 & 2 & 0 & 0 & 0 & 0\\
% 	0 & 0 & 0 & 2 & 1 & 0 & 0\\
% 	0 & 0 & 0 & 0 & 2 & 0 & 0\\
% 	0 & 0 & 0 & 0 & 0 & 3 & 0\\
% 	0 & 0 & 0 & 0 & 0 & 0 & 3
% 	\end{smatrix}
% 		\left(\begin{array}{@{}c@{}c@{}c@{}c@{}}
% 		\begin{array}{@{\ }ccc@{\ }|}
% 			2 & 1 & 0\\
% 			0 & 2 & 1\\
% 			0 & 0 & 2\\\hline
% 		\end{array}&&&\\
% 		&\begin{array}{|@{\ }cc@{\ }|}\hline
% 			2 & 1\\
% 			0 & 2\\\hline
% 		\end{array}&&\\
% 		&&\begin{array}{|@{\ }c@{\ }|}\hline
% 			3\\\hline
% 		\end{array}&\\
% 		&&&\begin{array}{|@{\ }c@{\ }}\hline
% 			3
% 		\end{array}
% 		\end{array}\right)
		\]
	\begin{enumerate}
	  \item Find the characteristic polynomial of $\rT$.
	  \item Find the dot diagram for each eigenvalue.
	  \item For each eigenvalue find the smallest $k_j$ such that $K_{\lambda_j}=\cN(\rT-\lambda_j\I)^{k_j}$.
	\end{enumerate}
	  
	\item For each matrix $A$ find a Jordan canonical form and an invertible $Q$ such that $A=QJQ^{-1}$.
	\begin{enumerate}
	  \item $A=\begin{pmatrix}
	  -3 & 3 & -2\\
		-7 & 6 & -3\\
		1 & -1 & 2
	  \end{pmatrix}$\qquad
	  (b)\ \ $A=\begin{pmatrix}
	  0 & 1 & -1\\
		-4 & 4 & -2\\
		-2 & 1 & 1
	  \end{pmatrix}$\qquad
	  (c)\ \ $A=$\scalebox{0.9}{$\begin{pmatrix}
	  0 & -3 & 1 & 2\\
		-2 & 1 & -1 & 2\\
		-2 & 1 & -1 & 2\\
		-2 & -3 & 1 & 4
	  \end{pmatrix}$}
	\end{enumerate}
	
	\item For each linear operator $\rT$, find a Jordan canonical form $J$ and basis $\beta$:
	\begin{enumerate}
	  \item $\rT(f)=f'$ on $\Span_\R\{e^t,te^t,t^2e^{t},e^{2t}\}$
	  
	  \item $\rT\bigl(f(x)\bigr)=xf''(x)$ on $P_3(\R)$
	  
	  \item $\rT(f)=af_x+bf_y$ on $\Span_\R\{1,x,y,x^2,y^2,xy\}$. How does your answer depend on $a,b$?
	\end{enumerate}
	
	\item (Generalized Eigenvector Method for ODEs)\quad Let $A\in M_n(\R)$ have an eigenvalue $\lambda$ and suppose $\beta_{\vv_0}=\{\vv_{k-1},\ldots,\vv_1,\vv_0\}$ is a cycle of generalized eigenvectors for this eigenvalue. Show that
	\[\vx(t):=e^{\lambda t}\sum_{j=0}^{k-1}b_j(t)\vv_j\quad\text{satisfies}\quad\vx'(t)=A\vx\iff \begin{cases}
	b_0'(t)=0,\text{ and}\\
	b'_j(t)=b_{j-1}(t)\text{ when $j\ge 1$}
	\end{cases}\]
	Use this method to solve the system of differential equations
	\[\vx'=\scalebox{0.9}{$\begin{pmatrix}
	3&1&0&0\\
	0&3&1&0\\
	0&0&3&0\\
	0&0&0&2
	\end{pmatrix}\vx$}\]
\end{enumerate}
\end{exercises}

%\setcounter{subsection}{3}
 
\subsection{The Rational Canonical Form (non-examinable)}

We finish the course with a very quick discussion of what can be done when the characteristic polynomial of a linear map does not split. In such a situation, we may assume that
\[p(t)=(-1)^n\bigl(\phi_1(t)\bigr)^{m_1}\cdots\bigl(\phi_k(t)\bigr)^{m_k}\tag{$\ast$}\]
where each $\phi_j(t)$ is an \emph{irreducible monic polynomial} over the field.

\begin{example}{}{ratcon1}
The following matrix has characteristic equation $p(t)=(t^2+1)^2(3-t)$
\[A=\begin{smatrix}
0 & -1 & 0 & 0 & 0 &\\
1 & 0 & 0 & 0 & 0 &\\
0 & 0 & 0 & -1 & 0\\
0 & 0 & 1 & 0 & 0\\
0 & 0 & 0 & 0 & 3 
\end{smatrix}\in M_5(\R)\]
This doesn't split over $\R$ since $t^2+1=0$ has no real roots. It is, however, diagonalizable over $\C$.
\end{example}

A couple of basic facts from algebra:
\begin{itemize}\itemsep0pt
  \item Every polynomial splits over $\C$: every $A\in M_n(\C)$ therefore has a Jordan form.
  \item Every polynomial over $\R$ factorizes into linear or irreducible quadratic factors.
  %\item Over other fields, there can be higher degree irreducible polynomials.
\end{itemize}

The question is how to deal with non-linear irreducible factors in the characteristic polynomial.

\begin{defn}{}{}
The monic polynomial $t^k+a_{k-1}t^{k-1}+\cdots+a_0$ has \emph{companion matrix}
\[\scalebox{0.82}{$\begin{pmatrix}
0 & 0 & 0 & \cdots & 0 & -a_0\\
1 & 0 & 0 & & 0 & -a_1\\
0 & 1 & 0 & & 0 & -a_2\\[-3pt]
\vdots & & & \ddots & & \vdots\\[2pt]
0 & 0 & 0 & & 0 & -a_{k-2}\\
0 & 0 & 0 & \cdots & 1 & -a_{k-1}
\end{pmatrix}$}
\tag{when $k=1$, this is the $1\times 1$ matrix $(-a_0)$}\]
If $\rT\in\cL(V)$ has characteristic polynomial ($\ast$), then a \emph{rational canonical basis} is a basis for which
\[[\rT]_\beta=
\scalebox{0.8}{$\begin{pmatrix}
C_1&O&\cdots&O\\
O&C_2&&O\\
\vdots&&\ddots&\vdots\\
O&O&\cdots&C_r
\end{pmatrix}$}
\]
where each $C_j$ is a companion matrix of some $(\phi_j(t))^{s_j}$ where $s_j\le m_j$. We call $[\rT]_\beta$ a \emph{rational canonical form} of $\rT$.
\end{defn}

We state the main result without proof:
\begin{thm}{}{}
A rational canonical basis exists for \emph{any} linear operator $\rT$ on a finite-dimensional vector space $V$. The canonical form is unique up to ordering of companion matrices.
\end{thm}


\begin{example*}{\ref{ex:ratcon1} cont}{}
The matrix $A$ is \emph{already in rational canonical form}: the standard basis is rational canonical with three companion blocks,
\[C_1=C_2=\begin{smatrix}
0&-1\\
1&0
\end{smatrix},\quad C_3=(3)\]
\end{example*}

\begin{example}{}{}
Let $A=\begin{smatrix}
4&-3\\2&2
\end{smatrix}\in M_2(\R)$. Its characteristic polynomial
\[p(t)=t^2-6t+14=(t-3)^2+5\]
doesn't split over $\R$ and so it has no eigenvalues. Instead simply pick a vector, $\vx=\stwovec 10$ (say), define $\vy=A\vx=\stwovec 42$, let $\beta=\{\vx,\vy\}$ and observe that
\[[\rL_A]_\beta=\begin{pmatrix}
0&-14\\1&6
\end{pmatrix}\]
is a rational canonical form. Indeed this works for \emph{any} $\vx\neq\V0$: if $\beta:=\{\vx,A\vx\}$, then Cayley--Hamilton forces
\[A^2\vx=(6A-14I)\vx=-14\vx+6A\vx\implies  [\rL_A]_\beta=\begin{pmatrix}
0&-14\\1&6
\end{pmatrix}\]
whence $\beta$ is a rational canonical basis and the form $[\rL_A]_\beta$ is independent of $\vx$!
\end{example}

A systematic approach to finding rational canonical forms is similar to that for Jordan forms: for each irreducible divisor of $p(t)$, the subspace $K_\phi=\cN\bigl(\phi(\rT)\bigr)^m$ plays a role analogous to a generalized eigenspace; indeed $K_\lambda=K_\phi$ for the \emph{linear} irreducible factor $\phi(t)=\lambda-t$!\smallbreak
We finish with two examples; hopefully the approach is intuitive, even without theoretical justification.

\begin{examples}{}{}
If the characteristic polynomial of $\rT\in\cL(\R^4)$ is
\[p(t)=(\phi(t))^2=(t^2-2t+3)^2=t^4-4t^3+10t^2-12t+9\]
then there are two possible rational canonical forms; here is an example of each.
\begin{enumerate}
  \item If $A=\begin{smatrix}
  0&-15&0&-9\\
  2&2&-3&0\\
  0&-9&0&-6\\
  -3&0&5&2
  \end{smatrix}$, then $\phi(A)=O$ is the zero matrix, whence $\cN(\phi(A))=\R^4$. Since $\phi(t)$ isn't the full characteristic polynomial, we expect there to be \emph{two} independent cycles of length two in the canonical basis. Start with something simple as a guess:
  \[\vx_1=\sfourvec 1000\implies \vx_2=A\vx_1=\sfourvec 020{-3}\implies A\vx_2=\sfourvec{-3}40{-6}=-3\vx_1+2\vx_2\]
  Now make another choice that isn't in the span of $\{\vx_1,\vx_2\}$:
  \[\vx_3=\sfourvec 0010\implies \vx_4=A\vx_3=\sfourvec 0{-3}05\implies A\vx_4=\sfourvec 0{-6}{-3}{10}=-3\vx_3+2\vx_4\]
  We therefore have a rational canonical basis $\beta=\{\vx_1,\vx_2,\vx_3,\vx_4\}$ and
  \[A=
  \begin{smatrix}
  1&0&0&0\\
  0&2&0&-3\\
  0&0&1&0\\
  0&-3&0&5
  \end{smatrix}
  \begin{smatrix}
  0&-3&0&0\\
  1&2&0&0\\
  0&0&0&-3\\
  0&0&1&2
  \end{smatrix}
  \begin{smatrix}
  1&0&0&0\\
  0&2&0&-3\\
  0&0&1&0\\
  0&-3&0&5
  \end{smatrix}^{-1}\]
  Over $\C$, this example is diagonalizable. Indeed each of the $2\times 2$ companion matrices is diagonalizable over $\C$. 
  
  \item Let $B=\begin{smatrix}
  0&0&2&1\\
  1&1&-1&-1\\
  0&1&-2&-16\\
  0&0&1&5
  \end{smatrix}$. This time
  \[\phi(B)=B^2-2B+3I=\begin{smatrix}
  3&2&-7&-29\\
  -1&1&4&13\\
  1&-3&-6&-17\\
  0&1&1&2
  \end{smatrix}\implies \cN(\phi(B))=\Span\left\{\sfourvec 3{-1}10,\sfourvec{11}{-2}01\right\}\]
  Anything not in this span will suffice as a generator for a single cycle of length four: e.g.,
  \begin{gather*}
  \vx_1=\sfourvec 1000,\quad \vx_2=B\vx_1=\sfourvec 0100,\quad \vx_3=B\vx_2=\sfourvec 0110,\quad \vx_4=B\vx_3= \sfourvec 20{-1}1\\
  B\vx_4=\sfourvec{-1}2{-14}4=-9\sfourvec 1000+12\sfourvec 0100-10\sfourvec 0110+4\sfourvec 20{-1}1
  \end{gather*}
  We therefore have a rational canonical basis $\beta=\{\vx_1,\vx_2,\vx_3,\vx_4\}$ and
  \[B=
  \begin{smatrix}
  1&0&0&2\\
  0&1&1&0\\
  0&0&1&-1\\
  0&0&0&1
  \end{smatrix}
  \begin{smatrix}
  0&0&0&-9\\
  1&0&0&12\\
  0&1&0&-10\\
  0&0&1&4
  \end{smatrix}
  \begin{smatrix}
  1&0&0&2\\
  0&1&1&0\\
  0&0&1&-1\\
  0&0&0&1
  \end{smatrix}^{-1}\]
  In contrast to the first example, $B$ isn't diagonalizable over $\C$. It has Jordan form \smash{$J=\begin{smatrix}
  \lambda&1&0&0\\
  0&\lambda&0&0\\
  0&0&\cl\lambda&1\\
  0&0&0&\cl\lambda
  \end{smatrix}$}
  where $\lambda=1+i\sqrt 2$. 
\end{enumerate}
\end{examples}


% 
%  \subsection{The Minimal Polynomial}
%  
%  Wed week 10
